\documentclass[11pt]{article}
\usepackage[T1]{fontenc}
\usepackage[paperwidth=180mm,paperheight=210mm,margin=13mm]{geometry}
\usepackage{lmodern,amsmath,amssymb,mathtools,microtype}
\pagestyle{empty}
\begin{document}
\noindent\small THE CLANKERS\hfill14 SEPTEMBER 2026
\vspace{3mm}
\begin{center}
{\LARGE\bfseries Retaining the zero is not\\[2mm] the same as proving a hit}\par
\vspace{3mm}
{\large Prime-power boundary comparison for Erd\H{o}s--Straus}
\end{center}
\noindent For every original labelled divisor state, retain
\[
 R=4a-p,\quad D_\alpha=\frac R{\gcd(R,g_\alpha)},\qquad
 g_E=4u+1,\quad g_M=u+a.
\]
A hit means $D_\alpha=1$. With all prime-power defects retained,
\[
 J_\alpha=\mathbb Q[z_\ell]_\ell/
       (z_\ell^{d_{\alpha\ell}+1})_\ell,\qquad
 D_\alpha=\prod_\ell\ell^{d_{\alpha\ell}},
\]
\[
 \boxed{e_\alpha\ \xmapsto{\ s\ }\ z^{d_\alpha}e_\alpha
 \ \xmapsto{\ \epsilon\ }\ \mathbf1_{D_\alpha=1}e_\alpha.}
\]
Every failure still has a supported top line. The \emph{comparison}, not the
existence of that line, detects the successful original state.

\vspace{3mm}\hrule\vspace{3mm}
\noindent Put $x_\alpha=1/D_\alpha\in(0,1/3]\cup\{1\}$. With the original
arithmetic masses, canonical polynomial bounds give
\[
 \boxed{L_d\le H=\#\{\alpha:D_\alpha=1\}\le U_d.}
\]
\[
 U_d=\min_{\substack{\deg P\le d\\P(1)=1}}\sum_\alpha P(x_\alpha)^2,
 \qquad
 L_d=\max_{\substack{\deg P\le d\\P(1)=1}}
 \frac32\sum_\alpha(x_\alpha-1/3)P(x_\alpha)^2.
\]
\noindent Both extrema have exact rational normal equations, including singular
cases. Split representations carry their inverse-fibre weights.

\vspace{2mm}
\noindent\textbf{Example:}\quad $p=2521$, $a=636$, $R=23$.
\[
 \mu=3\delta_1+87\delta_{1/23},\qquad
 P(x)=\frac{23x-1}{22},\qquad L_1=U_1=3.
\]
\noindent\small Five matching split pairs, four residue-cohomology classes,
and three original divisor states remain three different objects.

\vspace{3mm}\hrule\vspace{2mm}
\noindent\footnotesize
Full proofs and exact verifier accompany the preprint.
The canonical upper minimum is classical Christoffel theory; the source
boundary and metric interfaces retain their SplitZero attribution.
No universal ES theorem, new prime coverage, or historical-priority claim.
\end{document}
